\section{The Axiom and the Two Constraints}
\label{sec:constraints}

The model has one axiom and two governing constraints. Every later result is forced, derived, or constructed from these three statements.

\begin{axiom}[Experience]
To exist is to experience or to be experienced. There is no inert, non-experiential existent.
\end{axiom}

This is the single substantive assumption. It makes experience not a property of some things but the meaning of being anything at all. Its two clauses, ``to experience'' and ``to be experienced,'' are not decoration. They become the two directions of every relation in the model (Section~\ref{sec:foundation}).

\begin{principle}[Discreteness]
The base is finite and discrete. There are finitely many distinctions at any moment, a finite tone alphabet, and discrete time. The continuum, meaning real numbers, smooth fields, and differential equations, is permitted only as a large-scale, emergent approximation, never as a foundation.
\end{principle}

\begin{principle}[Monism]
The only entity is the vibe. There is no separate link, field, force, particle, or spacetime substrate. The only primitive relation is the note. Anything that seems to be carried by a link, a field, or a particle must be a note or a structure of vibes.
\end{principle}

These two principles are filters on every move the model makes. Discreteness forbids smuggling in a continuum. Monism forbids smuggling in a second kind of thing. Together they force the whole construction to be built from one finite kind of stuff and one relation.

Monism has a corollary worth stating on its own, because it is easy to violate by habit.

\begin{principle}[No Hidden State]
A vibe carries no data beyond its tone. It has no private registers, timestamps, counters, stored coordinates, or attached metadata. Anything that looks like per-vibe bookkeeping is either emergent from the structure of notes or realized as a structure of additional vibes.
\end{principle}

In practice this means a vibe does not store a beat number or a clock, because the order of events lives in the connectivity of notes (Section~\ref{sec:time}). It does not store its address, because its position is its place in the note-structure (Section~\ref{sec:geometry}). It does not store an attention weight as private data, because attention is the tone of the note between two vibes (Section~\ref{sec:dynamics}). The test for any proposed quantity is simple. If it is not a tone, it must be either emergent from notes or itself made of vibes. Otherwise it is a forbidden new primitive.

A word on the word ``link.'' We keep it as convenient shorthand. A link is never a separate entity. It means a note (the bare relation) or a structure of vibes (a relation that has become a vibe in its own right). This is made precise in Section~\ref{sec:foundation}.
